% !PRINTED 537
\weberpage{537}

Eliminiren wir aus diesen linearen Gleichungen die $\omega_\nu$, so
können wir die Gleichung $n^{\text{ten}}$ Grades für $\lambda$ in Determinanten\-
form darstellen, und das Product der Wurzeln dieser Gleichung,
also die Norm von $\lambda$, erhalten wir als die Determinante aus
den $n^2$ Grössen
\[
\sum^{s} a_{s,\nu}\,t_{s,r}
\]
(\S.\,137, 8.). Diese Determinante lässt sich aber nach dem Multi\-
plicationssatze der Determinanten (s. Bd. I, \S.\,27) zerlegen, und
giebt, wenn
\[
A = \Sigma \pm a_{1,1}a_{2,2}\ldots a_{n,n}, \quad T = \Sigma \pm t_{1,1}t_{2,2}\ldots t_{n,n}
\]
gesetzt wird,
\weq{8}{N(\lambda) = A\,T.}

Hierin ist $T$ eine ganze Function der Variablen $t_\nu$ mit ganzen
rationalen Zahlenco\"efficienten, von der wir nun noch nachweisen
werden, dass sie primitiv ist.

Nehmen wir an, im Theiler von $T$ gehe irgend eine natür\-
liche Primzahl $p$ auf. Dann können wir $n$ ganze rationale
Formen $y_1, y_2, \ldots, y_n$ der Variablen $t$ so bestimmen, dass die
$n$ Summen
\weq{9}{u_s = t_{s,1}y_1 + t_{s,2}y_2 + \cdots + t_{s,n}y_n}
alle durch $p$ theilbar sind, ohne dass alle $y_1, y_2, \ldots, y_n$ durch $p$
theilbar sind.

Die Richtigkeit dieser Behauptung folgt aus elementaren
Determinantensätzen (Bd. I, zweiter Abschnitt).

Wenn nämlich die $t_{s,r}$ alle durch $p$ theilbar sind, so können
wir die $y_r$ ganz beliebig z.\,B. gleich 1 annehmen.

Andernfalls nehmen wir an, dass ausser der Determinante $T$
auch alle $m$-reihigen Unterdeterminanten durch $p$ theilbar seien
$(m \leqq n)$,
und dass unter den $(m-1)$-reihigen Unterdeter\-
minanten wenigstens eine nicht durch $p$ theilbar sei. Ist dann
etwa unter den $(m-1)$-reihigen Determinanten der Matrix
\[
\begin{array}{cccc}
t_{1,1}, & t_{1,2}, & \ldots, & t_{1,m} \\
\multicolumn{4}{c}{\cdots\cdots\cdots\cdots\cdots\cdots\cdots\cdots\cdots\cdots} \\
t_{m-1,1}, & t_{m-1,2}, & \ldots, & t_{m-1,m}
\end{array}
\]
eine durch $p$ nicht theilbar, so setzen wir für $y_1, y_2, \ldots, y_m$ eben
diese $(m-1)$-reihigen Determinanten und nehmen $y_{m+1}, \ldots, y_n$
$= 0$ an. Diese $y$ genügen dann der gestellten Forderung.
